\section{Proof of Theorem~\ref{thm:rigidity}}\label{app:rigidity-proof}
We separate the proof into three steps because continuity alone does not permit us to differentiate the encoder. A fourth transmitter first turns equality of encoded triple sums into equality of both moments. We then use topology to obtain local dependence on those moments and regularization to justify differentiation. The resulting derivative identities determine the encoder coefficients.

\begin{lemma}
\label{lem:collision}
Let $K\ge4$, $I=[-1,1]$, $\phi:I\to\R$, and $\psi:I\to\R$, where $\psi$ is strictly convex or strictly concave. Suppose $f_{\phi,\psi}:I^K\to\R$ in \eqref{eq:target} is identifiable through $\enc:I\to\R^2$. Then, for every $\mathbf a,\mathbf b\in(-1,1)^3$,
\begin{align*}
 &\sum_{j=1}^3\enc(a_j)=\sum_{j=1}^3\enc(b_j)\\
 &\quad\Longrightarrow\quad
 \begin{pmatrix}\sum_{j=1}^3 a_j\\\sum_{j=1}^3\phi(a_j)\end{pmatrix}
 =\begin{pmatrix}\sum_{j=1}^3 b_j\\\sum_{j=1}^3\phi(b_j)\end{pmatrix}.
\end{align*}
\end{lemma}
\begin{proof}
Let the triple statistics be $(S,T)$ and $(\widetilde S,\widetilde T)$. Complete both triples with the same fourth input $w$ and $K-4$ zeros. The identifiability condition gives
$$
 T-\widetilde T
 =K\left[\psi\!\left(\frac{S+w}{K}\right)
 -\psi\!\left(\frac{\widetilde S+w}{K}\right)\right].
$$
Use $w=0$ and $w=1/2$. With $a=S/K$, $\widetilde a=\widetilde S/K$, and $\delta=1/(2K)>0$, subtraction yields
$$
 \psi(a+\delta)-\psi(a)
 =\psi(\widetilde a+\delta)-\psi(\widetilde a).
$$
All four arguments lie in $(-1,1)$. If $a<\widetilde a$, the points $a+\delta$ and $\widetilde a$ lie strictly between $a$ and $\widetilde a+\delta$ with complementary convex weights. Strict convexity therefore gives
$$
 \psi(a+\delta)+\psi(\widetilde a)
 <\psi(a)+\psi(\widetilde a+\delta),
$$
contradicting the equal increments. Strict concavity reverses this inequality, and exchanging the triples excludes $a>\widetilde a$. Thus $S=\widetilde S$, and the same identifiability condition at $w=0$ gives $T=\widetilde T$. This proof also permits endpoint inputs and does not require differentiability of $\psi$.
\end{proof}

The collision lemma ensures that the received sum determines the moments. We must also show that it cannot vary locally while those moments remain fixed. The next lemma supplies this converse dependence using invariance of domain: a continuous injection from an open subset of $\R^2$ into $\R^2$ has open image~\cite[Theorem 2B.3]{hatcher}.

\begin{lemma}
\label{lem:topology}
Let $B\subset\R^2$ be open, let $J\subset\R$ be an interval, and let $R:B\times J\to\R^2$ be continuous. Suppose there is a map $q:R(B\times J)\to B$ such that $q(R(p,t))=p$ for every $(p,t)\in B\times J$. Then $R(p,t)$ is independent of $t$ for each fixed $p$.
\end{lemma}
\begin{proof}
For fixed $t_0$, the map $p\mapsto R(p,t_0)$ is a continuous injection, so $R(B,t_0)$ is open. Fix $p$. For $t$ sufficiently close to $t_0$, continuity gives $R(p,t)\in R(B,t_0)$. Thus $R(p,t)=R(p',t_0)$ for some $p'\in B$. Applying $q$ yields $p'=p$, so $R(p,t)=R(p,t_0)$. The function $t\mapsto R(p,t)$ is locally constant, and connectedness of $J$ makes it constant. This argument does not require differentiability or continuity of $q$.
\end{proof}

The final lemma adapts restricted Cauchy--Pexider rigidity~\cite[Theorem 1.6 and Sec. 5]{charalambides2013} to a one-way collision implication and a nonaffine local preprocessing function $\phi$ of class $C^1$ that may contain affine subintervals.

\begin{lemma}
\label{lem:affine}
Let $\phi:I\to\R$ be nonaffine and belong to $C(I)\cap C^1((-1,1))$, and let $\enc:I\to\R^2$ be continuous. Suppose equality of encoded sums for any two interior triples implies equality of their $S$ and $T$ statistics. Then $\enc$ has form \eqref{eq:affine} with $c,u,v\in\R^2$ and linearly independent $u,v$.
\end{lemma}
\begin{proof}
\emph{Step 1: Local factorization through the moments.}
Choose interior $y_0,z_0$ with $\phi'(y_0)\ne\phi'(z_0)$; such a pair exists because $\phi$ is nonaffine. At any $(x_0,y_0,z_0)$, the map
$$
 (x,y,z)\longmapsto
 \bigl(x,\;x+y+z,\;\phi(x)+\phi(y)+\phi(z)\bigr)
$$
has Jacobian determinant $\phi'(z_0)-\phi'(y_0)\ne0$ and is a local $C^1$ diffeomorphism. Restrict its image to a product $J\times B$, where $x_0\in J$ and $B\subset\R^2$ is open. In these coordinates, the triple received sum is a continuous $R(p,x)$. The collision hypothesis makes the assignment $q(R(p,x))=p$ well defined even when distinct points of the chart have the same received sum. Lemma~\ref{lem:topology} shows that $R$ is independent of $x$. In particular, $H(p)=R(p,x_0)$ is continuous. On a sufficiently small source box whose image lies in $J\times B$, we have
\begin{equation}
 \enc(x)+\enc(y)+\enc(z)
 =H\bigl(x+y+z,\phi(x)+\phi(y)+\phi(z)\bigr)
 \label{eq:localfactor}
\end{equation}
for a continuous $H$ defined on $B$.

\emph{Step 2: Automatic interior differentiability.}
The pair map $(y,z)\mapsto w=(y+z,\phi(y)+\phi(z))$ is a local $C^1$ diffeomorphism near $(y_0,z_0)$. Restrict this map to the pair-coordinate part of the source box from Step~1. Choose a smooth compactly supported density $\rho(w)$ of integral one inside its open image, and denote its inverse by $(y(w),z(w))$. Integrating \eqref{eq:localfactor} with respect to $\rho(w)\,dw$ gives, for $x$ in the common source interval,
$$
 \enc(x)=\int H\bigl((x,\phi(x))+w\bigr)\rho(w)\,dw-C_0,
$$
where $C_0=\int[\enc(y(w))+\enc(z(w))]\rho(w)\,dw$ is independent of $x$. Choose the source interval $X$ and support of $\rho$ small enough that $\{(x,\phi(x))+w:x\in\overline X,\ w\in\operatorname{supp}\rho\}$ is compactly contained in $B$. Choose a smooth compactly supported cutoff in $B$ equal to one near this set, multiply $H$ by it, and extend the product by zero. The integral is then a smooth convolution as a function of the two-dimensional argument. Its composition with the $C^1$ curve $x\mapsto(x,\phi(x))$ is $C^1$. Since $x_0$ was arbitrary, $\enc$ is $C^1$ throughout the interval interior. No regularity of the original decoder was assumed.

\emph{Step 3: Affine dependence on $s$ and $\phi(s)$.}
Consider a real component $e$ of the now differentiable encoder. Equation~\eqref{eq:localfactor} makes $e(x)+e(y)+e(z)$ constant on each local moment fiber. The two-moment map has rank two whenever $\phi'(y)\ne\phi'(z)$, so the implicit-function theorem identifies its differential kernel with the tangent space of that fiber. The differential of the encoded sum annihilates this kernel. Hence
$$
 (e'(x),e'(y),e'(z))
 \in\operatorname{span}\{(1,1,1),(\phi'(x),\phi'(y),\phi'(z))\}.
$$
Fix $y=y_0,z=z_0$. Their entries give
$$
\beta=\frac{e'(y_0)-e'(z_0)}{\phi'(y_0)-\phi'(z_0)},
\qquad \alpha=e'(y_0)-\beta\phi'(y_0).
$$
These constants do not depend on $x$. Since Step~1 applies at every $(x,y_0,z_0)$, we obtain $e'(x)=\alpha+\beta\phi'(x)$ throughout the connected interval $(-1,1)$. Integration, followed by endpoint continuity, proves \eqref{eq:affine} componentwise.

Finally, the triple moment map has an open image locally near $(x_0,y_0,z_0)$. If $u,v$ were dependent, the linear map $(S,T)\mapsto uS+vT$ could not be injective on that open set. This would contradict the collision hypothesis. Thus $u,v$ are independent.
\end{proof}

\begin{proof}[Proof of Theorem~\ref{thm:rigidity}]
Necessity follows from Lemmas~\ref{lem:collision} and~\ref{lem:affine}. For sufficiency, the receiver solves the invertible real system $r_\enc-Kc=uS+vT$ and evaluates $T/K-\psi(S/K)$.
\end{proof}
